% \section{Future Work}
\label{sec:future_work}

\paragraph{Future Work.}
Several directions remain for future work.
\textbf{(1)} Our method can be extended to account for auxiliary factors that influence LLM-as-a-judge evaluations beyond the true label $\vZ$.
While our current analysis assumes that the LLM evaluation $\hat{\vZ}$ depends only on the true label $\vZ$,
in practice it may also be affected by additional nuisance factors $\xi$, such as response length or other stylistic attributes.
This more realistic setting can be accommodated by allowing both the confusion matrix
$\Pr_{\sQ}(\hat{\vZ} \mid \vZ,\xi)$ and the LLM evaluation $\E_{\sP}[\hat{\vZ} \mid \xi]$ to depend on such factors.
For example, one may estimate separate confusion matrices across different strata of $\xi$
(e.g., short vs.\ long responses) and perform bias adjustment within each group,
followed by aggregation across strata.
Formally, this corresponds to correcting $\E_{\sP}[\hat{\vZ} \mid \xi]$
using $\Pr_{\sQ}(\hat{\vZ} \mid \vZ,\xi)$ at the stratum level and averaging over $\xi$.
\textbf{(2)} The proposed method can be extended to a multinomial setting by generalizing $\big( \Pr_{\sQ}(\hat{\vZ} \mid \vZ) \big)^{-1} \E_{\sP}[\hat{\vZ}]$ in Section~\ref{sec:other:estimator} to a multivariate formulation, where $\vZ$ and $\hat{\vZ}$ are represented as probability distributions over multiple response categories. Such an extension would increase the number of parameters, and additional structural assumptions, such as constraints on the confusion matrix $\Pr_{\sQ}(\hat{\vZ} \mid \vZ)$, may be required to ensure stable estimation.
\textbf{(3)} A conformal prediction framework could be incorporated to provide sample-specific uncertainty quantification by constructing statistically valid prediction regions around each point estimate.
